\documentclass[11pt]{article}
\usepackage[margin=1in]{geometry}
\usepackage{graphicx}
\usepackage{booktabs}
\usepackage{siunitx}
\usepackage{amsmath}
\usepackage{hyperref}

\title{GNN Power-Flow Reconstruction with Droop/Deadband Inputs\\
Case24 IEEE RTS, 10{,}000 Scenarios}
\author{}
\date{}

\begin{document}
\maketitle

\begin{abstract}
This report summarizes the final Graph Neural Network (GNN) performance for power-flow feature reconstruction on the Case24 IEEE RTS dataset with 10{,}000 operating scenarios, including droop and deadband control parameters as node inputs. The model uses a transformer-based message-passing architecture and is trained for 100 epochs. Final evaluation on 1{,}000 held-out test scenarios (24{,}000 bus-level samples) shows strong accuracy across all targets: \(R^2\) values of 0.9987 (Pd), 0.9403 (Qd), 0.9996 (Pg), 0.9740 (Qg), 0.9951 (Vm), and 0.99999 (Va). A post-processing fix was applied so that unmasked variables are copied from known target values at inference, consistent with the masked-reconstruction objective. This removed systematic voltage-magnitude artifacts (notably at bus 10) and significantly improved all reported error metrics.
\end{abstract}

\section{Results}
\subsection{Setup and Test Split}
The final run uses:
\begin{itemize}
    \item Dataset: \texttt{case24\_ieee\_rts\_s10000\_droop\_mprange\_mqrange}
    \item Model: \texttt{GNN\_TransformerConv}
    \item Input dimension: 13 (including \texttt{mp\_droop}, \texttt{mq\_droop}, \texttt{frequency\_deadband}, \texttt{voltage\_deadband})
    \item Training epochs: 100
    \item Test set: 1{,}000 scenarios (24 buses each, 24{,}000 rows)
\end{itemize}

The merged test table is saved as \texttt{test\_true\_pred\_merged.csv}. Aggregate metrics are in \texttt{summary\_rmse\_mae\_bias.csv}.

\subsection{Main Accuracy Metrics}
\begin{table}[h]
\centering
\caption{Final test-set metrics (after clean inference fix).}
\begin{tabular}{lrrrr}
\toprule
Target & RMSE & MAE & Bias & \(R^2\) \\
\midrule
Pd & 3.5674 & 1.4480 & -0.2615 & 0.9987 \\
Qd & 4.9348 & 1.8344 & -0.0858 & 0.9403 \\
Pg & 3.6147 & 1.5396 & -0.5994 & 0.9996 \\
Qg & 7.7429 & 3.1218 & -0.2863 & 0.9740 \\
Vm & 0.0015 & 0.0008 & -0.00008 & 0.9951 \\
Va & 0.0311 & 0.0162 & -0.0010 & 0.99999 \\
\bottomrule
\end{tabular}
\end{table}

\subsection{Improvement vs. Pre-fix 100-Epoch Inference}
The file \texttt{comparison\_vs\_v2\_before\_fix.csv} compares the same trained model before and after the inference-output fix. RMSE reductions are:
\begin{itemize}
    \item Pd: -31.50\%
    \item Qd: -28.25\%
    \item Pg: -70.35\%
    \item Qg: -30.09\%
    \item Vm: -98.31\%
    \item Va: -33.22\%
\end{itemize}

\subsection{Figures}
Figures produced in this run are:
\begin{itemize}
    \item \texttt{true\_vs\_pred\_scatter.png}
    \item \texttt{error\_histograms.png}
    \item \texttt{true\_vs\_pred\_by\_bus\_one\_test\_scenario.png}
\end{itemize}

\begin{figure}[h]
    \centering
    \includegraphics[width=0.95\textwidth]{true_vs_pred_scatter.png}
    \caption{True vs. predicted values for all six targets on the test set.}
\end{figure}

\begin{figure}[h]
    \centering
    \includegraphics[width=0.95\textwidth]{error_histograms.png}
    \caption{Prediction error distributions (\( \hat{y}-y \)) for all targets.}
\end{figure}

\begin{figure}[h]
    \centering
    \includegraphics[width=0.95\textwidth]{true_vs_pred_by_bus_one_test_scenario.png}
    \caption{Per-bus true vs. predicted trajectories for one representative test scenario.}
\end{figure}

\section{Discussion}
\subsection{What Worked}
Increasing training from 5 epochs to 100 epochs substantially improved generalization. The model learns active and reactive power quantities accurately and provides very high-fidelity voltage-angle estimation. Including droop/deadband parameters in node features is compatible with the current architecture and contributes to stable performance over a large scenario set.

\subsection{Root Cause of Earlier Vm Outliers}
Earlier plots showed large Vm outliers (especially around bus 10). The root cause was not data corruption; it was an inference-consistency issue in masked reconstruction. During prediction, unmasked entries were left as raw model outputs, even though training losses emphasize masked targets and use physics residuals with unmasked values clamped to targets. Consequently, some unmasked outputs drifted without being directly constrained.

\subsection{Why the Fix is Correct}
The implemented fix enforces task consistency: at test/predict time, unmasked entries are copied from known target values, while masked entries remain model predictions. This aligns reported outputs with the masked-reconstruction objective and removes artificial outlier modes. The post-fix metrics confirm this, with the largest gain on Vm and clear improvements on all other targets.

\subsection{Limitations and Next Steps}
Although overall accuracy is strong, Qd and Qg remain the hardest targets relative to Pd/Pg/Va. Recommended next steps are: (i) sensitivity analysis over mask strategy (\texttt{rnd} vs \texttt{pf}), (ii) ablation on droop/deadband features, and (iii) hyperparameter sweeps for reactive-power-focused loss weighting.

\section*{Artifacts}
All outputs are saved in:
\begin{itemize}
    \item \texttt{results/gnn\_case24\_droop\_10000\_v3\_100ep\_fixed}
\end{itemize}
Key files:
\begin{itemize}
    \item \texttt{summary\_rmse\_mae\_bias.csv}
    \item \texttt{comparison\_vs\_v2\_before\_fix.csv}
    \item \texttt{test\_true\_pred\_merged.csv}
    \item \texttt{run\_info.txt}
\end{itemize}

\end{document}
